% Adapter-local capacity and layer-prefix allocation, part of section 3.6.
\paragraph{Control Peer capacity each window.}
For a target DP $p$ in window $t$, its Adapter grants
\begin{equation}
 B_{p,t}=[\min(B_{\mathrm{cap}},\mathrm{pool}_{p,t}
              -\mathrm{watermark}_{p,t}-\mathrm{headroom}_{p,t})]^+,
 \label{eq:peer-budget}
\end{equation}
where $B_{\mathrm{cap}}$ defaults to 4\,GiB. The pool includes bytes
already occupied by Peer replicas; $B_{p,t}$ is a total target, not
additional free space. Foreground demand, admission headroom, and
transfer/compute estimates are refreshed each window.

Let $A_r(k)$ contain the KV Master's hot sets in layers $1\ldots k$
and all selector state needed by those layers, with $A_r(0)=\varnothing$.
Given the useful prefix limit $L_r^\star(t)$ derived below, the Adapter
processes $Q_p$ in the Master's priority order. Starting with $P_s=\varnothing$,
it chooses
\begin{equation}
 \begin{aligned}
 L_r(t)=\max\bigl\{&k\in\{0,\ldots,L_r^\star(t)\}:\\
                  &\mathrm{bytes}\bigl(P_s\cup A_r(k)\bigr)\le B_{p,t}\bigr\},
 \end{aligned}
 \label{eq:resident-prefix}
\end{equation}
then adds $A_r(L_r(t))$ to $P_s$. Shared physical objects count once.
Only complete per-layer hot sets are admitted, and $L_r(t)=0$ preserves
the DRAM path. A request can receive fewer than $L_r^\star(t)$ layers;
unused space can serve later requests. Capacity or profile changes can
increase or decrease $L_r(t)$ at the next window. This greedy allocation
has no global optimality claim. After shared-object accounting, it scans
at most $N$ layer prefixes per request using marginal byte counts.

Within a window, these are stable targets. Emergency foreground pressure
can shrink them early: the Adapter removes later layers first, starting
with requests of lowest admission priority. Copying, committed, and
reclaiming bytes remain physically charged until release completes;
lowering $B_{p,t}$ or $L_r(t)$ does not immediately free addresses.
New copies wait for sufficient physical headroom, and readers retain
their pins until use completes. Growth copies only newly required
objects; every removal preserves the required DRAM recovery path.
